\chapter{A Concrete Naming Certificate for $\LocatedBisectionLimitUp$}
\label{ch:concrete-instances-located-bisection-limit-namecert}
\origin{human}

Following Bishop and Bridges on interval bisection and constructive real approximation, the $\LocatedBisectionLimitUp$ packet names the finite certificate surface by which a retained located bisection route is read as a terminal $\RealUp$ seal. It sits over the source located-interval bisection packet of \autoref{ch:concrete-instances-located-interval-bisection-namecert}, the bisection-to-Cauchy handoff of \autoref{ch:concrete-instances-bisection-cauchy-namecert}, the dyadic tolerance ledger of \autoref{ch:concrete-instances-dyadicratcore-namecert}, the finite observation interface of \autoref{ch:concrete-instances-streamname-namecert}, the regular rational readback of \autoref{ch:concrete-instances-regseqrat-namecert}, and the terminal real-name surface of \autoref{ch:concrete-instances-real-namecert}. The packet records a finite route from located bisection data to a Real seal; it does not select an infinite branch, identify quotient streams, assert trichotomy, invoke countable choice, or claim ambient $\RealUp$ completeness.

\begin{definition}[Located bisection limit carrier]
\label{def:located-bisection-limit-carrier}
A $\LocatedBisectionLimitUp$ carrier is a finite $\BHist$ packet
$$
L=(I,B,D,S,R,E,H,C,P,N)
$$
with the following displayed rows:
\begin{enumerate}
\item $I$ is the located interval source row, read through the $\LocatedIntervalBisectionUp$ surface and carrying the displayed endpoint, radius, locator, and retained-window data for the requested bisection depth.
\item $B$ is the retained bisection branch ledger. It records the finite left-or-right branch prefix, the retained child intervals, and the source agreement with $I$.
\item $D$ is the $\DyadicRatCoreUp$ midpoint and diameter budget row. It stores the displayed midpoint cells, dyadic width bounds, and tolerance comparisons used by the retained branch ledger.
\item $S$ is the $\StreamNameUp$ window row listing the finite bisection stages and endpoint approximants inspected by a downstream request.
\item $R$ is the $\RegSeqRatUp$ readback row formed from the retained dyadic centers and the diameter budgets displayed in $D$.
\item $E$ is the terminal $\RealUp$ seal. It is exposed only after $I,B,D,S,R$ have been displayed and replayed.
\item $H$ records componentwise $\hsame$ transport for the source interval, retained branch ledger, dyadic midpoint and diameter budget, StreamName windows, RegSeqRat readback, terminal Real seal, provenance, and local naming rows.
\item $C$ records displayed $\Cont$ replay from a finite approximation request through the located interval source, retained branch ledger, dyadic budget, StreamName windows, RegSeqRat readback, and terminal Real seal.
\item $P$ is the $\Pkg$ provenance over $\LocatedBisectionLimitUp$, $\LocatedIntervalBisectionUp$, $\BisectionCauchyUp$, $\DyadicRatCoreUp$, $\StreamNameUp$, $\RegSeqRatUp$, $\RealUp$, $\BHist$, $\hsame$, $\Cont$, and $\NameCert$.
\item $N$ is the local $\NameCert_{\LocatedBisectionLimitUp}$ row.
\end{enumerate}
The classifier compares accepted packets only through $I,B,D,S,R,E,H,C,P,N$. It has no coordinate for an infinite branch object, quotient stream equality, host real equality, trichotomy, countable choice, selected representatives, or ambient $\RealUp$ completeness.
\end{definition}

\begin{theorem}[Located bisection limit NameCert obligations]
\label{thm:located-bisection-limit-namecert-obligations}
For every accepted $\LocatedBisectionLimitUp$ carrier
$$
L=(I,B,D,S,R,E,H,C,P,N),
$$
the local $\NameCert_{\LocatedBisectionLimitUp}$ surface consists exactly of the following finite obligation rows:
\begin{enumerate}
\item carrier admission for the located interval source, retained bisection branch ledger, dyadic midpoint and diameter budget, finite stream windows, regular rational readback, terminal Real seal, componentwise transport, replay, provenance, and local naming rows;
\item source agreement, requiring the retained branch ledger $B$ to read its initial interval and locator data from $I$;
\item dyadic diameter control, requiring every retained child interval in $B$ to display midpoint and width information through $D$ before any readback row is exposed;
\item finite-window extraction, requiring $S$ to list only bisection stages and endpoint approximants already displayed by $I,B,D$;
\item regular-readback handoff, requiring $R$ to read retained centers under the dyadic error budgets supplied by $D$ before the terminal seal $E$ is exposed;
\item transport, replay, provenance, and naming stability through $H,C,P,N$;
\item non-escape, excluding selected infinite branches, quotient stream equality, host real equality, trichotomy, countable choice, selected representatives, and ambient $\RealUp$ completeness from the local certificate.
\end{enumerate}
No NameCert obligation is stored outside $I,B,D,S,R,E,H,C,P,N$.
\end{theorem}
\begin{proof}
Carrier admission is projection from \autoref{def:located-bisection-limit-carrier}. The packet displays the located source interval, retained branch ledger, dyadic budget, stream windows, regular readback, terminal seal, and structural rows.

Source agreement is inherited from the located-interval bisection route of \autoref{ch:concrete-instances-located-interval-bisection-namecert}. The branch ledger has no independent interval source; it can only retain child intervals below the displayed row $I$.

Dyadic diameter control follows the bisection-to-Cauchy route of \autoref{ch:concrete-instances-bisection-cauchy-namecert} and the tolerance ledger of \autoref{ch:concrete-instances-dyadicratcore-namecert}. Each retained stage displays a midpoint and width before a downstream window may inspect it.

The finite-window and regular-readback obligations are supplied by \autoref{ch:concrete-instances-streamname-namecert} and \autoref{ch:concrete-instances-regseqrat-namecert}. The window row selects only displayed stages, and the readback row turns those stages into regular rational data with the same dyadic budget.

Finally the terminal handoff is the $\RealUp$ seal of \autoref{ch:concrete-instances-real-namecert}, transported and replayed through $H,C,P,N$. Since the carrier has no row for the excluded branch, quotient, choice, trichotomy, host-equality, or completeness objects, the listed rows exhaust the local NameCert surface.
\end{proof}

\begin{theorem}[Located bisection limit Real-seal handoff]
\label{thm:located-bisection-limit-real-seal-handoff}
Every $\RealUp$-facing consumer of an accepted $\LocatedBisectionLimitUp$ carrier must pass through the finite route
$$
I,B,D \longmapsto S,R \longmapsto E,H,C,P,N .
$$
The consumer receives located interval data, a retained bisection branch ledger, dyadic midpoint and diameter budgets, finite StreamName windows, RegSeqRat readback, componentwise $\hsame$ transport, displayed $\Cont$ replay, $\Pkg$ provenance, and a local $\NameCert$ row before the terminal Real seal can be consumed. The handoff exports no selected infinite branch, quotient stream equality, host real equality, trichotomy, countable choice, selected representative, or ambient $\RealUp$ completeness principle.
\end{theorem}
\begin{proof}
Start with the source and branch rows. The source row $I$ gives the located interval data, while $B$ displays the finite retained branch ledger below that source. This is the public bisection surface supplied by \autoref{ch:concrete-instances-located-interval-bisection-namecert}.

Next read the dyadic budget row $D$. The midpoint and diameter comparisons in $D$ are the data used by the bisection-Cauchy route of \autoref{ch:concrete-instances-bisection-cauchy-namecert} to control the retained centers.

Then the finite observation and regular readback rows $S,R$ expose only already displayed bisection stages. Their StreamName and RegSeqRat roles are the same finite-window and regular-rational roles described in \autoref{ch:concrete-instances-streamname-namecert} and \autoref{ch:concrete-instances-regseqrat-namecert}.

Finally $E,H,C,P,N$ seal and name the route. The terminal $\RealUp$ seal is reached only after the source, branch, dyadic, window, and readback rows have been replayed, as in \autoref{ch:concrete-instances-real-namecert}. No excluded branch, quotient, host-equality, trichotomy, choice, selected-representative, or completeness row is present, so no Real-facing consumer can bypass the displayed handoff.
\end{proof}

\closureat{LocatedBisectionLimitUp}{seedStr}

\begin{closurestatus}{\LocatedBisectionLimitUp}
  \constructivestory{A $\LocatedBisectionLimitUp$ packet is generated from a LocatedIntervalBisection source row, a retained bisection branch ledger, DyadicRatCore midpoint and diameter budgets, StreamName finite windows, RegSeqRat readback, a terminal Real seal, componentwise hsame transport, displayed Cont replay, Pkg provenance, and a local NameCert row.}
  \theoryclosure{\seedClosure}
  \scopeclosed{\autoref{def:located-bisection-limit-carrier}, \autoref{thm:located-bisection-limit-namecert-obligations}, and \autoref{thm:located-bisection-limit-real-seal-handoff} seed the finite located-bisection route from source interval and dyadic branch rows to regular Cauchy readback and terminal Real-seal consumers.}
  \formalstatus{\unformalizedV}
  \bridgestatus{none}
  \notclaimed{This chapter does not close $\RealUp$ completeness, selected infinite branches, quotient stream equality, trichotomy, countable choice, host real equality, selected representatives, Lean verification, or standard bridge checking.}
  \upgradepath{Obligation closure requires checked carrier admission, source agreement, dyadic diameter control, finite-window extraction, regular-readback handoff, Real-seal handoff, and non-escape for BEDC.Derived.LocatedBisectionLimitUp.}
\end{closurestatus}
